% !TeX root = ../概率论与数理统计笔记.tex
% 由主文档通过 \input 载入；请编译主文档。
\setcounter{chapter}{15} % 保留原书章号。
\chapter{线性回归与预测}
\label{chap:ps-regression}

本章用线性模型研究响应变量如何随解释变量变化，并量化估计与预测误差。
先修内容为条件期望、协方差、估计和检验；多元部分需要矩阵秩、正交投影与二次型。

\section{回归问题与模型假设}
\label{sec:ps-regression-model}
定义\emph{响应变量}（response variable）、\emph{解释变量}（explanatory variable）与\emph{线性回归}（linear regression）。
先采用固定设计的一元模型 $Y_j=\beta_0+\beta_1x_j+\varepsilon_j$：$x_j$ 为给定设计值，$\beta_0,\beta_1$ 为未知系数，$\varepsilon_j$ 为随机误差。
分别说明零均值、同方差、不相关或独立、正态性在不同结论中的作用；随机设计情形改用给定解释变量的条件假设。

\section{最小二乘估计与几何解释}
\label{sec:ps-least-squares}
定义\emph{最小二乘法}（least squares），明确优化变量为回归系数、目标为观测残差平方和。
在设计值不全相同时推导一元回归解；多元模型以设计矩阵列满秩保证系数解唯一。
用正交投影解释拟合值与残差，区分线性代数优化事实与需要随机误差假设的统计结论。

\section{估计量的分布与高斯--马尔可夫定理}
\label{sec:ps-gauss-markov}
推导回归系数估计量的期望、协方差与误差方差估计。
在固定满秩设计、误差零均值及协方差为标量单位阵等条件下，陈述高斯--马尔可夫定理（Gauss--Markov theorem）的最佳线性无偏性。
精确 $t$ 与 $F$ 推断另加正态误差假设；最佳线性无偏不等于在所有估计量中最优。

\section{回归系数推断与预测区间}
\label{sec:ps-regression-inference}
构造斜率和一般线性系数的置信区间及检验，介绍整体回归显著性的 $F$ 检验。
区分给定解释变量处平均响应的区间与新观测的\emph{预测区间}（prediction interval），说明后者还包含新误差的波动。
预测新观测时明确其与训练误差的关系，并说明外推超出数据范围时依赖更强的模型可信度。

\section{诊断、多元回归与模型局限}
\label{sec:ps-regression-diagnostics}
介绍残差图、异常观测、异方差、共线性与影响点；多元模型说明“控制其他变量”的系数解释所依赖的模型。
定义决定系数并限定其解释，说明较高拟合程度不能单独证明预测性能或因果关系。
模型选择、稳健标准误与训练验证划分作为后续拓展，说明其解决的问题而不把软件输出当作自动结论。

\section{例题、习题与学习检查}
规划一元直线拟合、给定设计点的均值与预测区间、多元设计矩阵秩不足的例子。
推导题分别检查最小化、无偏性和区间覆盖所需条件；应用题要求把变量单位、估计系数与不确定性一起解释。
学习检查区分误差与残差、相关与回归、描述关系与因果解释。






